\chapter{Using \LaTeX~for scientific documents}
\label{chap:latex}

\section{What's \LaTeX?} 

In your research, you will produce papers, reports and -- very 
importantly -- your thesis. These documents can be written using a 
WYSIWYG (What You See Is What You Get) editor (e.g., Word). However, an 
alternative especially suited for scientific publications is \LaTeX. In 
\LaTeX, the document is simply a text file ({\tt .tex}). Text 
formatting is using markups (like HTML). The file is then ``compiled'' 
(like source code of a programming language) into a file -- typically 
{\tt .pdf}. 

\section{Why \LaTeX?}

\begin{compactitem}
  \item The input is a small, portable text file
  \item \LaTeX~compilers are freely available for all OS'
  \item Exactly the same result on any computer (not true for Word)
  \item \LaTeX~produces beautiful, professional looking docs (e.g., like 
this one!)
  \item Mathematical formulas (esp complex ones) are easy to write
  \item \LaTeX~is very stable -- current version basically same since 1994! 
(9 major versions of MS Word since 1994 -- with compatibility issues)
    \item \LaTeX~is free!  
    \item Many journals provide \LaTeX~templates, making formatting quicker 
    \item Bibliographies are a breeze and work with Mendeley and Zotero
    \item Plenty of online support available -- your question has probably 
already been answered
    \item You can integrate \LaTeX~into a workflow to auto-generate lengthy 
and complex documents (like your thesis).  
\end{compactitem}

\begin{center}	
	\includegraphics[width=.4\textwidth]{miktex.png}
\end{center}

\subsection{Limitations of \LaTeX}

\begin{compactitem}\itemsep10pt
  \item It has a steeper learning curve.
  
  \item Can be difficult to manage revisions with multiple authors -- 
especially if they don't use \LaTeX! (I have a dark secret)
    
  \item Typesetting tables can be a bit complex.

    \item Images and floats don't jump like Word, but if you don't use the 
right package, they can be difficult to place where you want! 

  \end{compactitem}

\section{Installing \LaTeX}
 
\begin{lstlisting} 
sudo apt-get install texlive-full texlive-fonts-recommended 
latex-beamer texlive-pictures texlive-latex-extra imagemagick       
\end{lstlisting}

We will use a text editor in this lecture, but you can use one of a 
number of WYSIWYG frontends (e.g., {\tt Lyx}, {\tt TeXmacs}), as well 
as GUI's ({\tt texmaker}, {\tt Gummi}, {\tt TeXShop}, etc). Overleaf (\url{https://www.overleaf.com/})
(now merged with Sharelatex) is also very good (and works with git), especially for 
collaborating with non \LaTeX~-ers (your university may have a blanket 
license for the pro version).  

\section{A first \LaTeX~example}

\begin{compactitem}[$\quad\star$]\itemsep4pt{}
  \item Open geany and type the following in a file
	{\tt Week1/Code/FirstExample.tex}:
\end{compactitem}

\lstinputlisting{Practicals/Code/FirstExample.tex}

Now, let's get a citation for Einstein's paper:

\begin{compactitem}[$\quad\star$]\itemsep4pt{}
\item In Google Scholar, go to ``settings'' (upper right corner) and choose
BibTeX as bibliography manager. 
  \item Now type ``does the energy of a body einstein 1905''
  \item The paper should be the one on the top. 
  \item Click `` Import into BibTeX'' should show the following text, that you
will save in the file {\tt FirstBiblio.bib} (in the same directory as
{\tt FirstExample.tex})
{\tiny \lstinputlisting{Practicals/Code/FirstBiblio.bib}}
  \end{compactitem}

Now we can create a {\tt .pdf} of the article. 
\begin{compactitem}[$\quad\star$]\itemsep4pt{}
  \item In the terminal type (are you the right directory?!):
  \end{compactitem}

\begin{lstlisting}
$ pdflatex FirstExample.tex
$ pdflatex FirstExample.tex
$ bibtex FirstExample
$ pdflatex FirstExample.tex
$ pdflatex FirstExample.tex
\end{lstlisting}

This should produce the file {\tt FirstExample.pdf}:
    
\begin{center}
\setlength\fboxsep{0pt}
\setlength\fboxrule{0.5pt}
\fbox{\includegraphics[width = 0.5\linewidth]{FirstExample.pdf}}
\end{center}

\subsection{A bash script to compile \LaTeX}

You can of course write a useful little bash script to compile latex 
with bibtex! 

Type the following script and call it {\tt CompileLaTeX.sh} (you know 
where to put it!):

\lstinputlisting[language=bash]{Practicals/Code/CompileLaTeX.sh}

How do you run this script? The same as your previous bash scripts, so 
\begin{lstlisting}	
$ bash CompileLaTeX.sh FirstExample
\end{lstlisting}

Why have I not written the {\tt *.tex} extension of {\tt FirstExample} 
in the command above?  

\section{A brief \LaTeX~tour}

\begin{compactitem}\itemsep12pt
\item Spaces, new lines and special characters:

\begin{compactitem} \setlength{\itemindent}{-1em}\itemsep2pt
  \item Several spaces in your text editor are treated as one space in the
typeset document
  \item Several empty lines are treated as one empty line
  \item One empty line defines a new paragraph
  \item Some characters are ``special'':\# \$ \% \^{} \& \_ \{ \}
\~{} \textbackslash. 
  \item To type these special characters, you have to add a ``backslash'' in
front, e.g., \textbackslash\$ produces \$.
\end{compactitem}

\item Document structure:

\begin{compactitem} \setlength{\itemindent}{-1em}\itemsep2pt
  \item Each \LaTeX~command starts with \textbackslash (e.g., to get \LaTeX,
you need {\tt \textbackslash{LaTeX}}
  \item The first command is always {\tt \textbackslash{documentclass}} defining
the type of document (e.g., {\tt article, book, report, letter}).
  \item You can set several options. For example, to set size of text to 10
points and the letter paper size: {\tt
\textbackslash{documentclass[10pt,letterpaper]\{article\}}}.

\end{compactitem}

  \item After having declared the type of document, you can specify packages
    you want to use. The most useful are:

\begin{compactitem}
  \item{\tt \textbackslash{usepackage\{color\}}}: use colors for text in your
document.
  \item{\tt \textbackslash{usepackage\{amsmath,amssymb\}}}: American
Mathematical Society formats and commands for typesetting mathematics. 
  \item{\tt \textbackslash{usepackage\{fancyhdr\}}}: fancy headers and footers. 
  \item{\tt \textbackslash{usepackage\{graphicx\}}}: include figures in pdf, ps,
eps, gif and jpeg.
  \item{\tt \textbackslash{usepackage\{listings\}}}: typeset source code for
various programming languages.
  \item{\tt \textbackslash{usepackage\{rotating\}}}: rotate tables and figures.
  \item{\tt \textbackslash{usepackage\{lineno\}}}: line numbers.
\end{compactitem}

\item Once you select the packages, you can start your document with {\tt
\textbackslash{begin\{document\}}}, and end it with {\tt
\textbackslash{end\{document\}}}.

\end{compactitem}

\section{\LaTeX~templates}

There a lots of useful \LaTeX templates out there. As an example of 
structure of a document, take the article template provided by the 
journal PNAS:

{\tiny \lstinputlisting{Practicals/Code/PnasExample.tex}}

I have added some templates in the {\tt CMEEMasteRepo} that you should 
have a look and play around with

\section{Typesetting math}

There are two ways to display math 

 \begin{enumerate}

  \item First, one can produce inline mathematics (i.e., within the text).
    \item Second, one can produce stand-alone, numbered equations and formulae.
    
   \end{enumerate}
   
For inline math, the ``dollar'' sign flanks the math to be typeset. For 
example, the code:
\begin{lstlisting}
$\int_0^1 p^x (1-p)^y dp$
\end{lstlisting}

Becomes $\int_0^1 p^x (1-p)^y dp$

For numbered equations (almost always a great idea), \LaTeX~provides 
the {\tt equation} environment:

\begin{lstlisting}

\begin{equation}
\int_0^1 \left(\ln \left( \frac{1}{x} \right) 
\right)^y dx = y!
\end{equation}

\end{lstlisting}

Becomes:
\begin{equation}
\int_0^1 \left(\ln \left( \frac{1}{x} \right) \right)^y dx = y!
\end{equation}

\section{A few more tips}

The following tips might prove handy:

\begin{compactitem}
\item \LaTeX~has a full set of symbols and operators (plenty of lists online)
\item Long documents can be split into separate {\tt.tex} documents and
combined using {\tt input}
\item Long documents can be split into separate {\tt.tex} documents and
\item Figures can be included using the {\tt graphicx} package
\item You can use Mendeley to export and maintain {\tt .bib} files
\item You can redefine environments and commands in the preamble 
\end{compactitem}

\subsection[Practical]{Practicals}
  
	Test {\tt CompileLaTeX.sh} with {\tt FirstExample.tex} and 
	bring it under verson control under{\tt CMEECourseWork/Week1} in 
	your repository. Make sure that {\tt CompileLaTeX.sh} will work if 
	I ran it from my computer using {\tt FirstExample.tex} as an input.

\section[Practical]{Practicals wrap-up}

Make sure you have your {\tt Week 1} directory organized 
with {\tt Data}, {\tt Sandbox} and {\tt Code} with the necessary 
files and this week's (functional!) scripts in there. Every script 
should run without errors on my computer. This includes the five 
solutions (single-line commands you came up with) in {\tt 
UnixPrac1.txt}.

\begin{center}
	\it Commit and push everything by next Wednesday 5 PM.
\end{center} 

\section{Readings \& Resources}

\begin{itemize}

\item The Visual \LaTeX~ FAQ: sometimes it is difficult to describe what you
want to do!\\
\url{http://get-software.net/info/visualFAQ/visualFAQ.pdf}

\item Myriad online resources for \LaTeX, including:\\
\url{www.http://en.wikibooks.org/wiki/LaTeX/Introduction},\\
\url{www.ctan.org/tex-archive/info/lshort/english/}\\
\url{http://ftp.uni-erlangen.de/mirrors/CTAN/info/lshort/english/lshort.pdf}
\item Beautiful presentations in \LaTeX: 
\url{http://tug.org/pracjourn/2005-2/miller/miller.pdf}
\item Bibliographies in \LaTeX: \url{http://schneider.ncifcrf.gov/latex.html}

\end{itemize}


